%
% mode-b-call.tex
%
% Z specification of the Mode B voice call in voxd
% (src/punt_vox/voxd/voice_call/, PL-1371): the call lifecycle that
% call_start, call_stop and call_status drive, the playback ownership of
% AgentSpeaker, and the per-call tool result cache. The tests in
% tests/voice_call/ assert the named properties.
%

\documentclass[a4paper,10pt,fleqn]{article}
\usepackage[margin=1in]{geometry}
\usepackage{fuzz}

\begin{document}

\title{The Mode B Voice Call: A Z Specification}
\author{Punt Labs, for PL-1371}
\date{November 2026}
\maketitle

\section{Introduction}

voxd runs at most one Mode B call. \texttt{call\_start} creates it in
\emph{connecting}; the call task provisions the ElevenLabs agent and opens
the conversation, which makes it \emph{live}. A remote hang-up, a
\texttt{call\_stop} or a connect timeout ends it; any other error fails it.
While a call holds the output, it owns the playback mutex across replies;
an interruption ends the current reply only. Each tool call is keyed by its
canonical arguments; a durable key runs its operation at most once per call.

The model names six properties, and the tests carry the same names:

\begin{enumerate}
  \item[C1] \emph{One call}: a second \texttt{call\_start} is refused while a
    call is connecting or live (\texttt{test\_call\_barge\_in\_reask\_starts\_one\_worker}
    asserts it through the handlers).
  \item[C2] \emph{Stoppable at once}: \texttt{call\_stop} ends a call that is
    still provisioning or connecting
    (\texttt{test\_stop\_cancels\_a\_call\_still\_provisioning},
    \texttt{test\_stop\_cancels\_a\_call\_still\_connecting}).
  \item[C3] \emph{Terminal states stay}: an ended or failed call never
    becomes live again.
  \item[C4] \emph{Output held across a cut}: an interruption ends the reply
    but not the call's hold on playback
    (\texttt{test\_the\_call\_keeps\_the\_output\_across\_a\_cut}).
  \item[C5] \emph{Output released at the end}: no ended or failed call holds
    playback.
  \item[C6] \emph{Durable once}: a re-asked durable tool call replays and
    never runs its operation a second time
    (\texttt{test\_reask\_during\_a\_slow\_start\_starts\_one\_worker}).
\end{enumerate}

\section{Types}

\begin{zed}
  [KEY] \\
  CALL ::= idle | connecting | live | ended | failed \\
  ZBOOL ::= ztrue | zfalse
\end{zed}

\texttt{idle} is the handlers' state before the first call.

\section{State}

\begin{schema}{Call}
  call : CALL \\
  holding, replying : ZBOOL \\
  durable, ran : \power KEY
\where
  call \in \{ idle, ended, failed \} \implies holding = zfalse \\
  replying = ztrue \implies holding = ztrue \\
  ran \subseteq durable
\end{schema}

The first predicate is C5. \texttt{durable} is the set of durable keys whose
operation has run or is in flight; \texttt{ran} those that finished.

\begin{schema}{InitCall}
  Call'
\where
  call' = idle \\
  holding' = zfalse \\
  replying' = zfalse \\
  durable' = \emptyset \\
  ran' = \emptyset
\end{schema}

\section{Lifecycle}

\begin{schema}{Start}
  \Delta Call
\where
  call \in \{ idle, ended, failed \} \\
  call' = connecting \\
  holding' = zfalse \\
  replying' = zfalse \\
  durable' = \emptyset \\
  ran' = \emptyset
\end{schema}

\texttt{Start} is refused from \emph{connecting} or \emph{live} (C1), and a
new call gets a fresh cache.

\begin{schema}{Connect}
  \Delta Call
\where
  call = connecting \\
  call' = live \\
  holding' = holding \\
  replying' = replying \\
  durable' = durable \\
  ran' = ran
\end{schema}

\begin{schema}{End}
  \Delta Call \\
  failure? : ZBOOL
\where
  call \in \{ connecting, live \} \\
  failure? = ztrue \implies call' = failed \\
  failure? = zfalse \implies call' = ended \\
  holding' = zfalse \\
  replying' = zfalse \\
  durable' = durable \\
  ran' = ran
\end{schema}

\texttt{End} covers hang-up, stop, timeout and failure. Its precondition
admits \emph{connecting}, which is C2; no operation leaves \emph{ended} or
\emph{failed} except a new \texttt{Start}, which is C3.

\section{Playback}

\begin{schema}{Speak}
  \Delta Call
\where
  call = live \\
  holding' = ztrue \\
  replying' = ztrue \\
  call' = call \\
  durable' = durable \\
  ran' = ran
\end{schema}

\begin{schema}{Interrupt}
  \Delta Call
\where
  call = live \\
  replying' = zfalse \\
  holding' = holding \\
  call' = call \\
  durable' = durable \\
  ran' = ran
\end{schema}

\texttt{Interrupt} leaves \texttt{holding} as it was: C4.

\section{Tool results}

\begin{schema}{RunDurable}
  \Delta Call \\
  k? : KEY
\where
  call = live \\
  k? \notin durable \\
  durable' = durable \cup \{ k? \} \\
  ran' = ran \\
  call' = call \\
  holding' = holding \\
  replying' = replying
\end{schema}

\begin{schema}{Replay}
  \Xi Call \\
  k? : KEY
\where
  call = live \\
  k? \in durable
\end{schema}

\begin{schema}{Finish}
  \Delta Call \\
  k? : KEY
\where
  k? \in durable \setminus ran \\
  ran' = ran \cup \{ k? \} \\
  durable' = durable \\
  call' = call \\
  holding' = holding \\
  replying' = replying
\end{schema}

A key enters \texttt{durable} only through \texttt{RunDurable}, whose
precondition is $k? \notin durable$; a re-ask finds it there and replays,
whether the first run finished or is still in flight under the shield: C6.

\end{document}
